\documentclass[aspectratio=169]{beamer}
\input{../theme/imsb_beamer.tex}

\title{Imaging-based single-cell spatial data: Xenium}
\author{Dr.\ Robin Khatri}
\institute{Institute of Medical Systems Bioinformatics, UKE}
\date{Day 2 \textendash{} 17.06.2026}

\begin{document}

\begin{frame}[plain]\titlepage\end{frame}

\begin{frame}{Recap and today}
  Yesterday: spot-based Visium, the standard workflow, and first spatial
  statistics.
  \vskip0.5em
  Today we move to \textbf{imaging-based} data at single-cell resolution:
  \begin{itemize}
    \item how the data are produced and why segmentation is the key step;
    \item quality control and evaluation metrics specific to this technology;
    \item a full single-cell spatial workflow, and the method choices along the way.
  \end{itemize}
\end{frame}

\begin{frame}{How imaging-based ST works}
  \begin{columns}[T]
  \begin{column}{0.5\textwidth}
  \begin{itemize}
    \item A targeted gene panel is detected \emph{in situ} by cyclic hybridisation
      and imaging of fluorescent probes.
    \item Each gene has a code; across rounds the code is read out and decoded to a
      transcript at an $(x,y)$ location.
    \item Transcripts are assigned to cells using a segmentation mask.
    \item Output: transcripts, a cell $\times$ gene matrix, cell metadata, and images.
  \end{itemize}
  \end{column}
  \begin{column}{0.48\textwidth}
  \centering
  \begin{tikzpicture}[scale=0.8]
    \foreach \i/\x in {1/0,2/1.3,3/2.6,4/3.9}{
      \node[draw=imsbblue,thick,minimum size=0.7cm] (r\i) at (\x,2.6) {};
      \node[lab,below=0pt of r\i]{R\i};}
    \fill[imsbblue!60] (0,2.6) circle (0.08);
    \fill[imsbblue!60] (2.6,2.6) circle (0.08);
    \fill[imsbblue!60] (3.9,2.6) circle (0.08);
    \node[box,below=8mm of r2,xshift=0.6cm] (code) {codeword\\1011};
    \node[box,below=5mm of code] (gene) {gene = CD3D};
    \draw[arr] (r2.south) -- (code.north);
    \draw[arr] (code) -- (gene);
  \end{tikzpicture}
  \end{column}
  \end{columns}
\end{frame}

\begin{frame}{Imaging-based platforms}
  \small
  \begin{center}
  \begin{tabular}{@{}lll@{}}
    \toprule
    Platform & Panel size & Note \\
    \midrule
    MERFISH / Vizgen & $\sim$300\textendash 1000 & error-robust barcodes \\
    seqFISH+ & $\sim$10{,}000 & many pseudocolours \\
    Xenium (10x) & 280\textendash 480, up to 5000 (Prime) & in situ, H\&E compatible \\
    CosMx (NanoString) & $\sim$1000\textendash 6000 & RNA and protein \\
    \bottomrule
  \end{tabular}
  \end{center}
  \vskip0.4em
  We use Xenium with the 377-gene Multi-Tissue and Cancer panel. The panel choice
  bounds what biology you can see, so it is part of the experimental design.
\end{frame}

\begin{frame}{Xenium outputs and how to read them}
  \begin{itemize}
    \item \texttt{cell\_feature\_matrix.h5}: cell $\times$ feature counts (genes and
      controls).
    \item \texttt{cells.parquet}: centroid, area, counts per cell.
    \item \texttt{transcripts.parquet}: every decoded transcript (large).
    \item \texttt{morphology\_*.ome.tif}: images (large).
  \end{itemize}
  \vskip0.3em
  \textbf{Two reading paths.} \texttt{spatialdata\_io.xenium} builds a full object
  with images and shapes (great for visualisation, heavy). For table-level analysis
  we read the matrix and cell table straight into AnnData. Same cells, a fraction
  of the memory: this is what keeps us inside 25 GB.
\end{frame}

\begin{frame}{Segmentation is the crux}
  \begin{columns}[T]
  \begin{column}{0.54\textwidth}
  Cells are defined by a mask, and everything downstream inherits its errors.
  \begin{itemize}
    \item In dense tissue, boundaries are ambiguous: transcripts are misassigned,
      and small or nucleus-less cells are missed.
    \item Symptoms: implausible co-expression, ``cells'' with too few transcripts,
      lineage markers bleeding across boundaries.
  \end{itemize}
  Treat cell-level results as estimates and sanity-check against the raw transcript
  image where it matters.
  \end{column}
  \begin{column}{0.44\textwidth}
  \centering
  \begin{tikzpicture}[scale=0.9]
    \draw[imsbdark,thick,dashed] (0,1.4) circle (1.05);
    \fill[imsbblue!40] (0,1.4) circle (0.45);
    \node[lab] at (0,1.4) {nucleus};
    \node[lab] at (0,0.05) {+ expansion};
    \foreach \x/\y in {0.9/1.9, 1.0/0.9, -0.95/1.7, -0.8/1.0, 0.2/2.3}{\fill[imsbgrey!70] (\x,\y) circle (0.07);}
    \node[lab,right=0.1cm] at (1.2,2.2){stray\\transcripts};
  \end{tikzpicture}
  \end{column}
  \end{columns}
\end{frame}

\begin{frame}{Segmentation methods}
  \small
  \begin{itemize}
    \item \textbf{Default (10x):} nuclear stain plus $\sim$15\,\textmu m expansion; a
      multimodal deep-learning kit uses membrane/boundary stains when available.
    \item \textbf{Image-based:} Cellpose (generalist nucleus/membrane segmentation).
    \item \textbf{Transcript-based:} Baysor and Proseg learn transcript
      colocalisation; Proseg uses a cellular Potts model.
    \item \textbf{Learned, reference-aware:} segger (graph neural network, transcript
      to cell link prediction, can use scRNA-seq).
  \end{itemize}
  \vskip0.3em
  Benchmarks show a tradeoff between transcript capture and cell purity that is
  tissue-dependent; there is no universal winner.
  \refline{Stringer et al., \emph{Nat Methods} 2021 (Cellpose); Petukhov et al. (Baysor); Jones et al.\ 2025 (Proseg); Heidari et al.\ 2025 (segger).}
\end{frame}

\begin{frame}{Why segmentation matters for your biology}
  \begin{itemize}
    \item Misassigned transcripts directly corrupt cell typing and, especially,
      cell\textendash cell communication (a ligand can ``jump'' to the wrong cell).
    \item Nucleus-anchored methods miss cells without a captured nucleus, biasing
      composition.
    \item Improving identity in poorly segmented, dense regions is an active and
      unsolved problem.
  \end{itemize}
  \figslot{the same field of view segmented by two methods, with cell counts and a marker that should not co-occur}{a segmentation-benchmark figure (e.g.\ segger or CRISP), or your own}
  \refline{Heidari et al.\ 2025 (segger); CRISP segmentation benchmark, bioRxiv 2026.}
\end{frame}

\begin{frame}{Control features}
  \begin{columns}[T]
  \begin{column}{0.56\textwidth}
  The panel ships with built-in negative controls:
  \begin{itemize}
    \item \textbf{Negative control probes}: target no real transcript; estimate
      non-specific binding.
    \item \textbf{Negative control codewords (blanks)}: valid codes assigned to no
      gene; estimate the decoding error rate.
  \end{itemize}
  A healthy run puts only a small fraction of counts on controls. A high control
  fraction in a cell or region is a warning.
  \end{column}
  \begin{column}{0.42\textwidth}
  \centering
  \begin{tikzpicture}[scale=0.9]
    \fill[imsbblue!55] (0,0) rectangle (0.6,3.2);
    \fill[imsbgrey!55] (1.1,0) rectangle (1.7,0.5);
    \node[lab,below] at (0.3,0){genes};
    \node[lab,below] at (1.4,0){controls};
    \node[lab,above] at (0.3,3.2){signal};
    \node[lab,above] at (1.4,0.5){noise};
  \end{tikzpicture}
  \end{column}
  \end{columns}
\end{frame}

\begin{frame}{Quality control and evaluation metrics}
  \begin{itemize}
    \item \textbf{Transcripts / cell} and \textbf{genes / cell}: too few means an
      empty or fragmented segment.
    \item \textbf{Control fraction / cell}: flags noise.
    \item \textbf{Cell area}: very small or very large segments are suspect.
    \item \textbf{Counts vs area}: density should be roughly consistent.
    \item \textbf{Panel sensitivity}: counts are far below scRNA-seq, so a gene's
      absence is weak evidence.
  \end{itemize}
  Read every metric in space, not only as a histogram.
\end{frame}

\begin{frame}{Cell typing on a small panel}
  \begin{itemize}
    \item Fewer genes make clusters noisier and rare types harder to separate.
    \item Use a curated marker set for the tissue, then refine by eye and check
      spatial plausibility.
    \item Reference label transfer (CellTypist, scANVI, \texttt{ingest}) scales when
      a matched scRNA-seq reference exists, but inherits any segmentation
      contamination.
    \item BANKSY can type cells and find domains together using the neighbourhood.
  \end{itemize}
  \refline{Dom\'inguez Conde et al., \emph{Science} 2022 (CellTypist); Singhal et al., \emph{Nat Genet} 2024 (BANKSY).}
\end{frame}

\begin{frame}{Choosing a spatial graph}
  \centering
  \begin{tikzpicture}[scale=0.8]
    \begin{scope}
      \foreach \x/\y in {0/0,0.8/0.6,0.3/1.2,1.2/1.4,1.6/0.4}{\fill[imsbblue!60] (\x,\y) circle (0.1);}
      \draw[line] (0,0)--(0.8,0.6); \draw[line] (0.8,0.6)--(0.3,1.2); \draw[line] (0.8,0.6)--(1.2,1.4); \draw[line] (0.8,0.6)--(1.6,0.4);
      \node[lab] at (0.7,-0.5){kNN};
    \end{scope}
    \begin{scope}[xshift=4cm]
      \foreach \x/\y in {0/0,0.8/0.6,0.3/1.2,1.2/1.4,1.6/0.4}{\fill[imsbblue!60] (\x,\y) circle (0.1);}
      \draw[line] (0,0)--(0.8,0.6)--(0.3,1.2)--(0,0); \draw[line] (0.8,0.6)--(1.2,1.4)--(0.3,1.2); \draw[line] (0.8,0.6)--(1.6,0.4)--(1.2,1.4);
      \node[lab] at (0.7,-0.5){Delaunay};
    \end{scope}
    \begin{scope}[xshift=8cm]
      \foreach \x/\y in {0/0,0.8/0.6,0.3/1.2,1.2/1.4,1.6/0.4}{\fill[imsbblue!60] (\x,\y) circle (0.1);}
      \draw[imsbgrey,dashed] (0.8,0.6) circle (0.9);
      \draw[line] (0.8,0.6)--(0,0); \draw[line] (0.8,0.6)--(0.3,1.2); \draw[line] (0.8,0.6)--(1.6,0.4);
      \node[lab] at (0.7,-0.5){radius};
    \end{scope}
  \end{tikzpicture}
  \vskip0.5em
  kNN fixes the neighbour count; Delaunay follows tissue geometry; radius fixes a
  physical distance. The choice changes enrichment results, so state it.
\end{frame}

\begin{frame}{Single-cell spatial statistics}
  \begin{itemize}
    \item \textbf{Neighbourhood enrichment}: which cell types co-locate vs avoid
      each other (permutation test on the graph).
    \item \textbf{Co-occurrence}: how that association changes with distance,
      separating tight contact from regional association.
    \item \textbf{Ripley's K / L}: clustering vs dispersion of a cell type across
      scales.
    \item \textbf{Spatially variable genes}: Moran's I over cells.
  \end{itemize}
  These are the building blocks for the organisation analysis on Day 3.
\end{frame}

\begin{frame}{Beyond one section}
  \begin{itemize}
    \item Multiple sections bring batch effects. Integration with Harmony or scVI
      aligns embeddings before joint clustering.
    \item Spatial-aware methods (e.g.\ BANKSY) can do batch-aware domain detection.
    \item Comparing conditions needs composition-aware tests, not just eyeballing
      proportions (Day 3).
  \end{itemize}
  \refline{Korsunsky et al., \emph{Nat Methods} 2019 (Harmony); Lopez et al., \emph{Nat Methods} 2018 (scVI).}
\end{frame}

\begin{frame}{Open problems on imaging-based data}
  \begin{itemize}
    \item Segmentation in dense tissue, and its propagation into typing and
      communication.
    \item Low detection of cytokines and chemokines, so immune states are hard to
      identify.
    \item Reliability of spatial QC metrics, which are still being standardised.
    \item Integrating images (H\&E) with the molecular layer for better region
      identity.
  \end{itemize}
\end{frame}

\begin{frame}{Further reading}
  \footnotesize
  \begin{itemize}
    \item Petukhov et al. Baysor: cell segmentation from transcripts.
    \item Jones et al.\ 2025 (Proseg); Heidari et al.\ 2025 (segger); Stringer et al., \emph{Nat Methods} 2021 (Cellpose).
    \item Singhal et al. BANKSY. \emph{Nat Genet} 2024.
    \item 10x Genomics analysis guides for Xenium segmentation and QC.
    \item Squidpy tutorials for Xenium: \texttt{squidpy.readthedocs.io}.
  \end{itemize}
\end{frame}

\begin{frame}{Hands-on: \texttt{day\_2/02\_tutorial.ipynb}}
  Dataset: a public Xenium non-diseased human kidney section ($\sim$97{,}560 cells).
  You will:
  \begin{itemize}
    \item read the section the lean way (matrix + cell table);
    \item separate control features and compute Xenium QC;
    \item filter, normalise, cluster, and annotate cell types with markers;
    \item plot cell types in space; build a spatial graph and run enrichment,
      co-occurrence, and Moran's I.
  \end{itemize}
  Exercises: set thresholds and map removed cells, refine annotation, compare kNN vs
  Delaunay graphs, and interpret spatially variable genes.
\end{frame}

\end{document}
